\section{Cira: contextual inference over residual memories}\label{sec:method}
\method{} separates three operations: identifying the active physical context, learning its dynamics correction, and transferring knowledge to contexts that share physical factors. The encoder and world-model predictor remain fixed throughout. Context inference selects what the robot already knows; predictable weighted updates determine what each memory learns next.

\subsection{Compact probabilistic residual memories}
Let $U\in\R^{d\times k}$ have orthonormal columns, with $k\ll d$, and let $\phi(z,a)\in\R^p$ be a fixed feature map. For compact notation, this section suppresses episode indices and orders transitions by deployment time. The projected residual is $v_t=U^\top r_t$, where $r_t=z_{t+1}-f(z_t,a_t)$. Memory $c$ predicts
\begin{equation}\label{eq:adapter}
 v_{t,i}\mid c,w_{c,i}\sim\mathcal N\!\left(w_{c,i}^{\top}\phi_t,\sigma_i^2(z_t,a_t)\right),
 \qquad i=1,\ldots,k,
\end{equation}
with $\phi_t=\phi(z_t,a_t)$. The rows $w_{c,i}^{\top}$ form $W_c\in\R^{k\times p}$; its estimated latent correction is $m_c(z,a)=U\widehat W_c\phi(z,a)$. Each row retains a Gaussian regression state. Independent rows and diagonal, state-dependent observation noise permit inexpensive updates while representing increased variability at contact. Decorrelation of training residuals motivates the diagonal approximation but does not establish independence at deployment.

The correction basis and feature map are fixed before the evaluation stream. Residual energy outside their span measures approximation error; a small number of changing physical parameters need not induce a low-rank latent discrepancy. The full predictive distribution also includes this uncorrected error. All safety scores are evaluated in the full latent space. The analysis in Section~\ref{sec:theory} states explicitly when exact linear realizability is required.

\subsection{A shared prior makes memories compositional}\label{sec:sharedprior}
For context $c=(c^1,\ldots,c^F)$, the prior couples memories that share factor values:
\begin{align}
 \operatorname{Cov}(w_{c,i},w_{c',i})&=k_{\mathcal C}(c,c')\Sigma_0,\label{eq:sharedprior}\\
 k_{\mathcal C}(c,c')&=\kappa_0^2+
 \sum_{j=1}^F\kappa_j^2\mathbf 1\{c^j=c'^j\}
 +\kappa_\times^2\mathbf 1\{c=c'\}.\label{eq:contextkernel}
\end{align}
Here $\Sigma_0\succ0$ is an input-weight covariance. This separable multi-output construction~\cite{alvarez2012kernels} contains a common correction, factor-specific components, and a combination-specific interaction. It gives one coherent prior over complete context memories, including tuples never observed together. The additive factor model corresponds to $\kappa_\times=0$; a positive interaction variance preserves uncertainty about unseen combinations rather than treating an additive estimate as exact.

For known combinations $\Omega$, joint Gaussian conditioning determines a new context's prior centre and covariance. With exactly known weights, its covariance coefficient is
\begin{equation}\label{eq:conditionalvariance}
 s_c=k_{\mathcal C}(c,c)-k_{\mathcal C}(c,\Omega)
 K_\Omega^{\dagger}k_{\mathcal C}(\Omega,c),
\end{equation}
where $K_\Omega$ is the context Gram matrix and $\dagger$ denotes the Moore--Penrose inverse. Uncertain memories require the joint Gaussian conditioning formula with their observation uncertainty. Replacing this joint state by separately fitted Gaussian memory summaries is a modular approximation: correlated summaries cannot be reused as independent observations without risking overconfidence. Proposition~\ref{prop:composition} concerns the exact conditional model.

Factor roles and cue regions are specified, such as a surface region and an object region; factor-value identities are learned from dynamics without labels. Independent variation, adequate excitation, and shared state--action support make these identities distinguishable. For two additive factors, the graph of observed combinations determines whether a new pair is identifiable. The kernel hyperparameters are fixed for the analyzed evaluation stream after calibration. Their uncertainty and misspecification are tested through sensitivity and coverage measurements. A reference context is imposed as a noise-free zero correction only when that correction is known; otherwise its uncertainty is retained.

\subsection{Inference before and after each transition}
The hypothesis library contains encountered contexts, unseen tuples of known factors, and novel-value hypotheses. The latter retain positive support when existing factor values cannot explain the dynamics. Joint novelty must remain possible when more than one factor changes. Memory allocation and hypothesis pruning are inference approximations whose missed-context rates are measured.

Let $q_t(c)$ denote the context probability immediately before observing $r_t$, and $\pi_t^+(c)$ the probability afterward. At an episode boundary,
\begin{equation}\label{eq:cueupdate}
 q_{\tau_e}(c)=\pi_{e,\tau_e^-}(c)
 \propto p(c\mid D_{<e})\prod_{j=1}^F p_\eta(y_e^j\mid c^j).
\end{equation}
The prior can encode frequency and persistence. During a fixed-context episode, $q_{t+1}=\pi_t^+$. A within-episode switching extension instead propagates $\pi_t^+$ through a transition model. The cue likelihood enters once per episode, before the first residual. After executing the action,
\begin{equation}\label{eq:dynamicsupdate}
 \pi_t^+(c)\propto q_t(c)\,p(v_t\mid z_t,a_t,c,\Theta_t),
\end{equation}
using memory predictions computed before learning from that transition. This posterior supports recognition, context switching, and novelty detection. Sustained support for a new-value hypothesis allocates a memory for future observations. An unseen combination of known values already has a prior prediction and can be selected before its first contact.

Cue models are updated from a parallel context inference pass that omits the episode's cue likelihood and uses dynamics evidence instead. Ambiguous assignments remain soft. This prevents direct training on the cue's own pseudo-labels, although policy-dependent data and errors in stored models can still propagate. Cue reversals and appearance-only changes test whether dynamics can correct these associations.

\subsection{Learning with probabilities fixed before the error}\label{sec:predictableupdates}
The inference posterior $\pi_t^+$ answers which context best explains the observed transition. Memory learning instead uses $q_t$, the probability fixed before that transition. For output row $i$, define
\begin{align}
 \omega_{c,i,t}&=q_t(c)/\sigma_i^2(z_t,a_t),\label{eq:weights}\\
 V_{c,i,t}&=P_{c,i}+\sum_{s<t}\omega_{c,i,s}\phi_s\phi_s^\top,\label{eq:precision}\\
 \widehat w_{c,i,t}&=V_{c,i,t}^{-1}\left(P_{c,i}\nu_{c,i,t}
                  +\sum_{s<t}\omega_{c,i,s}\phi_s v_{s,i}\right).\label{eq:weightedmean}
\end{align}
The prior centre $\nu_{c,i,t}$ transfers information through Eq.~\eqref{eq:sharedprior}. The positive-definite precision $P_{c,i}$ is fixed when the memory is instantiated, before its first update; a positive variance floor handles degenerate conditional priors. The centre may subsequently change using past information. The regression covariance $V_{c,i,t}^{-1}$ is distinct from the context posterior and from a guaranteed frequentist confidence region.

Only observations collected after instantiation enter these sums. Allocation-triggering residuals are not retrospectively selected into the new memory's regression statistics. Replaying a discovery window with assignments chosen from its own outcomes defines a separate ablation. Likewise, changing the precision requires a new analysis epoch with fresh sufficient statistics and a new failure-probability allocation.

This timing matters because a weight computed from the current error can select a truncated noise distribution. Predictable weights preserve the zero-mean noise condition used by Theorem~\ref{thm:confidence}. They do not eliminate contamination: after a context switch, probability may still favor the preceding condition. The theorem quantifies that cost in terms of accumulated incorrect assignment weight and links it to the context filter's log loss.

\subsection{Planning and decision-aware probing}\label{sec:planning}
For candidate committed action sequences $\mathcal A$, the planner evaluates context-specific rollouts and minimizes their posterior-weighted cost,
\begin{equation}\label{eq:planning}
 \widehat J_e(a)=\sum_c\pi_{e,\tau_e^-}(c)\widehat J_c(a).
\end{equation}
The analyzed planner rolls out each context's posterior-mean correction $m_c$, holding the context fixed over the horizon. Corollary~\ref{cor:decision} bounds this context-weighted objective over a fixed candidate set. A separate variant samples uncertain adapter parameters once per rollout and can optimize a risk-sensitive objective. For nonlinear dynamics and costs, that parameter-averaged objective requires its own approximation or uncertainty bound. No error observed during a committed strike can improve that strike.

A separate probe-enabled protocol permits an informative action before the consequential strike. With belief $b$ over context and memory parameters, let $V(b)=\min_{a\in\mathcal A}\E_b[J_\theta(a)]$. The informational value of probe $p$ is
\begin{equation}\label{eq:voi}
 \operatorname{VOI}(p)=V(b)-\E_{o\mid p,b}[V(b_o)],
\end{equation}
where $b_o$ is the Bayesian belief after its outcome. Among admissible probes, the robot selects the largest positive $\operatorname{VOI}(p)-\kappa(p)$, where $\kappa$ accounts for time and task cost. Proposition~\ref{prop:voi} identifies a useful boundary: reducing context uncertainty has zero decision value when all plausible models share an optimal action.

The computational approximation reuses the current candidate costs and updates only their context weights over sampled probe outcomes. It neglects state displacement, changes in feasible strikes, and within-memory learning. When uncertainty is dominated by a single memory, a local alternative targets the variance of the cost difference between the two best actions. For one row with covariance $\Sigma$ and cost-difference gradient $g$, its reduction is
\begin{equation}\label{eq:coptimal}
 \Delta_p=\frac{(g^\top\Sigma\phi_p)^2}
 {\phi_p^\top\Sigma\phi_p+\sigma_i^2(z_p,a_p)}.
\end{equation}
This is a linearized decision-variance criterion, not exact value of information. Both approximations are evaluated against mutual-information probing and no probing. Probe-enabled performance is reported separately from strict first-contact retrieval.

\subsection{Safety margins stored with each memory}\label{sec:safetymethod}
Let $h(z)\geq0$ define the latent safety condition, and let $\mathcal C_t^\delta$ be a nonempty set carrying at least $1-\delta$ posterior mass. For every plausible memory, the safety filter checks
\begin{equation}\label{eq:safety}
 h(\bar z'_{c,t})-L_h b_{c,t}\geq(1-\alpha)h(z_t),
 \qquad \bar z'_{c,t}=f(z_t,a_t)+m_{c,t}(z_t,a_t),
\end{equation}
where $L_h$ is a Lipschitz bound, $0<\alpha\leq1$, and $b_{c,t}$ is a remembered error margin. A linearized optimization can generate candidate actions, but its output requires exact constraint verification or a bounded linearization remainder to invoke Proposition~\ref{prop:safety}. A stopping or retraction action counts as a safe fallback only if its true safety condition is established.

After observing $z_{t+1}$, score each plausible memory against its own pre-update prediction:
\begin{equation}\label{eq:safetyscore}
 s_{c,t}=\|z_{t+1}-\bar z'_{c,t}\|_2.
\end{equation}
If any plausible memory covers the transition, assign the step to the most probable one satisfying $s_{c,t}\leq b_{c,t}$. Otherwise assign it to the most probable plausible memory and record a miss. Update only this memory's raw margin,
\begin{equation}\label{eq:marginupdate}
 b_c\leftarrow b_c+\gamma\left(\mathbf1\{s_{c,t}>b_{c,t}\}-\delta_s\right).
\end{equation}
This adaptive calibration recursion~\cite{gibbs2021adaptive} retains one margin per memory. The raw state is not clipped; the filter may conservatively replace it by $\max\{b_{c,t},0\}$. Predictive scores remain defined using the raw state for the calibration argument.

A familiar return recalls its previous margin. A wrong cue may omit the true context, but the full-state score still records the resulting mismatch. Under bounded scores and exact enforcement or a verified fallback, the calibration recursion bounds the cumulative number of one-step latent barrier failures for arbitrary cue sequences. This is a long-run accounting guarantee rather than protection against every first-contact failure. A physical guarantee additionally requires correspondence between $h$ and the physical constraint.

\subsection{Separating recall from parameter learning}\label{sec:decomposition}
Let $\pi^{\mathrm{pre}},\Theta^{\mathrm{pre}}$ be the adaptive state, including shared-prior centres, immediately before assimilating the current episode's cue, and $\pi^{\mathrm{post}}_s,\Theta^{\mathrm{post}}_s$ the state after $s$ post-contact observations. Thus $s=0$ on the updated side includes cue-based inference but no new transition. On a fixed held-out transition set $\mathcal Q_e$ for the current condition, evaluate four losses:
\begin{equation}
L_{ij}(s)=L_{\mathcal Q_e}(\pi^{(i)},\Theta^{(j)}),
\qquad i,j\in\{0,1\},
\end{equation}
where index $0$ means the pre-cue snapshot and index $1$ means the step-$s$ post-cue snapshot. Posterior weights are held fixed when parameters are exchanged; inference is not rerun. The total improvement $G(s)=L_{00}-L_{11}$ decomposes as
\begin{align}
R(s)&=\tfrac12\big[(L_{00}-L_{10})+(L_{01}-L_{11})\big],\label{eq:recall}\\
P(s)&=\tfrac12\big[(L_{00}-L_{01})+(L_{10}-L_{11})\big],\label{eq:proper}\\
G(s)&=R(s)+P(s).
\end{align}
These two-player Shapley contributions average the effect of each change across the two possible orders. On return episodes with an established library, $R$ measures the predictive contribution of recall (\emph{apparent learning}) and $P$ that of parameter updates (\emph{proper learning}). A method without memory selection has $R=0$ by construction. Either contribution can be negative, so both are reported with their signs.

Stable memory identities make the four counterfactuals comparable. When a new slot is created, its initial counterfactual uses a data-independent creation prior and zero initial selection mass; its fitted initialization cannot be inserted retrospectively into $\Theta^{\mathrm{pre}}$. In those windows, $R$ is labeled \emph{selection/creation}, since selecting an untrained slot is not recall of acquired knowledge. Creation events are reported separately, and the main recall analysis uses returns with fixed library membership. Held-out replay attributes predictive improvement conditional on the snapshots, not the causal contribution to closed-loop task success. A separate intervention that freezes all memory parameters during later returns directly tests whether retrieval alone retains the benefit.
